\subsection{Scientific Discovery Agents}
\label{sec:app-science}
Scientific-discovery agents aim to accelerate the entire life cycle for scientific research, from hypothesis generation through experimental execution, by coupling LLMs with domain-specific simulators, laboratory automation and up-to-date literature. These systems ground decision in verifiable processes while handling heterogeneous data, safety constraints and long-horizon goals.  

In this subsection, we begin with the foundational layer (Section~\ref{sec:app-science-planning}), which encompasses planning under scientific context, tool-augmented interaction with scientific resources, search and retrieval mechanisms including RAG-based systems and execution-time integration with laboratory hardware.
Building upon these capabilities, the self-evolving layer (Section~\ref{sec:app-science-memory}) introduces agentic memory, feedback and reflection, which enable scientific agents to refine hypotheses, adapt protocols and learn from experimental outcomes.
Finally, the collective layer (Section~\ref{sec:app-science-multi}) explores multi-agent collaboration, where agents coordinate roles, share intermediate knowledge and jointly reason toward complex scientific goals.



\subsubsection{Foundational agentic reasoning}
\label{sec:app-science-foundational}


\paragraph{Planning.}
\label{sec:app-science-planning}


Scientific agents utilize reasoning-enhanced planning ability to decompose a research goal into steps, decides which tool or simulator to call next, then revises the plan as evidence arrives. In short, the \textit{chain of thought} emerges from LLM reasoning that compiles instructions into rigorous executable plans~\cite{wei2022chain}.
% There are certain ways that LLM reasoning interacts with agentic planning ability. 
For example, ProtAgents~\cite{ghafarollahi2024protagents} materializes a planner agent that utilizes LLM reasoning capability to formulate a concrete plan for protein analysis and keep modifying it with feedback from another critic agent, and Eunomia~\cite{ansari2023agentbasedlearningmaterialsdatasets} uses ReAct-style~\cite{yao2023react} workflow to make in-context reasoning: after retrieving a top-$k$ evidence set, the backbone LLM quote a warranting sentence, and that citation drives the next action choice. Other examples include MatExpert~\cite{ding2024matexpertdecomposingmaterialsdiscovery}, which deploys a chain-of-thought LLM to author a stepwise transition pathway and then emits a structured crystal candidate from a feedback loop.

Planning can also act as reasoning constraint. For instance, Curie~\cite{kon2025curierigorousautomatedscientific} utilizes a rigor engine to align, setup and do reproducibility check within planning steps proposed by the Architect LLM. Thus, the Architect's free-form reasoning cannot advance unless these rigor gates are satisfied, which transforms planning into both a guide and a regulator of the reasoning process. In addition, a general purpose biomedical agent, Biomni~\cite{huang2025biomni} constrains its reasoning within a dynamically constructed biomedical action space of comprehensive tools,  software packages and databases, requiring each hypothesis to be operationalized as executable code.


% ------------------------------------------------------------------ %
\paragraph{Tool-Use.}
\label{sec:app-science-tool}



Tool use is an important part of the reasoning loop for scientific agents nowadays. Specifically, rather than following rigid rules,   these agents can decide which tool and when to call, how to fill parameters and verify or revise based on evidence. For example, SciAgent~\cite{ma2024sciagenttoolaugmentedlanguagemodels} formalizes \textit{tool-augmented reasoning} as a four-step procedure: planning, retrieval, tool-based action and execution. Agents are trained to decide when to call a tool, which one, and how to integrate it into solving scientific tasks. Through domain-specific tools, ChemCrow~\cite{bran2023chemcrowaugmentinglargelanguagemodels} chains various expert chemistry tools so intermediate calculations become premises in the next reasoning step, which enables end-to-end planning and autonomous syntheses. CACTUS~\cite{mcnaughton2024cactus} similarly grounds explanations in cheminformatics outputs, reducing reliance on free-form reasoning by language models alone. 

Other notable examples include ChemToolAgent~\cite{yu2024chemtoolagent} and CheMatAgent~\cite{wu2025chematagentenhancingllmschemistry}. In particular, ChemToolAgent~\cite{yu2024chemtoolagent} employs a ReAct-like~\cite{yao2023react} architecture with multiple specialized chemistry tools, allowing the LLM to choose and parameterize tool calls while CheMatAgent~\cite{wu2025chematagentenhancingllmschemistry} pushes further by learning tool use: it integrates over 100 chemistry/materials tools, curates a tool-specific benchmark, and uses Monte Carlo Tree Search with step-level fine-tuning to learn both which tool to pick and how to fill arguments.

For biomedical agents, TxAgent~\cite{gao2025txagentaiagenttherapeutic} scales therapeutic reasoning across 211 vetted tools and it carries out multi-step reasoning that reconciles drug labels, interactions, and patient context—turning clinical justification into an executable trace. On the other hand, AgentMD~\cite{jin2024agentmdempoweringlanguageagents} builds a two-stage tool memory: it first mines thousands of clinical calculators from literature (i.e. making tools), then selects and applies the right ones at inference (i.e. using tools), pinning predictions to concrete computations. Other recent systems~\cite{lála2023paperqaretrievalaugmentedgenerativeagent, skarlinski2024languageagentsachievesuperhuman,chiang2024llamp,zhang2024honeycombflexiblellmbasedagent,qu2025crisprgptagenticautomationgeneediting,gao2025pharmagentsbuildingvirtualpharma} reinforce similar design: co-design tool-use with reasoning so each claim is computable and auditable.



Another notable category of tool-use is the ability of agentic execution, which includes but not limited to run codes and simulate environments. Execution layers bridge high-level plans to physical infrastructure, which enables scientific agents to autonomously operate laboratory hardware, orchestrate simulation pipelines, and manage large-scale data workflows. Recent works such as Organa~\cite{darvish2025organaroboticassistantautomated} ties LLM reasoning to task-and-motion planning plus scheduling and perception, executing multi-step experiments with autonomous robots;  AtomAgents~\cite{ghafarollahi2024atomagentsalloydesigndiscovery} exemplifies the simulation side of execution: a physics-aware system that plans and runs atomistic workflows, coordinating tools for code execution, analysis, and hypothesis checking; and Chemist-x~\cite{chen2025chemistxlargelanguagemodelempowered} shows wet-lab execution beyond a digital-only scenario, where agents generate control scripts and drive an automated platform to validate conditions without human intervention. 

Several other platforms couple execution with optimization or team-based autonomy. For instance, SGA~\cite{ma2024llmsimulationbileveloptimizers} formalizes the workflow of \textit{LLM-as-proposer and simulator-as-optimizer}  while MatExpert~\cite{ding2024matexpertdecomposingmaterialsdiscovery} operates a \textit{retrieval, transition and generation} workflow for material discovery tasks, and CellAgent~\cite{xiao2024cellagentllmdrivenmultiagentframework} coordinates planner, executor and evaluator roles to run full single-cell analysis pipelines. 



\paragraph{Search and retrieval.}

\label{sec:app-science-search}



Beyond simple context stuffing, recent scientific agentic systems elevate retrieval into a deliberate reasoning step: agents decide when and what to fetch, and how to use the evidence before committing to a hypothesis. With retrieval ability, BioDiscoveryAgent~\cite{roohani2025biodiscoveryagentaiagentdesigning} pulls literature and interim assay results inside a closed loop so the model’s next gene-perturbation choices are conditioned on what was read and measured; while DrugAgent~\cite{inoue2025drugagentexplainabledrugrepurposing} coordinates knowledge graph queries, targeted literature search through web API and machine learning predictors. Its planner selects retrieval actions and then reconciles heterogeneous evidence into an explainable rationale. To facilitate scientific research, ARIA~\cite{ramirezmedina2025acceleratingscientificresearchmultillm} operationalizes a \textit{search, filter then synthesis} workflow as role-bound steps that carry citations forward, turning literature into actionable procedures. Similarly, AI Scientist-v2~\cite{yamada2025ai} employs an agentic tree-search framework in which the agent actively queries scientific literature database during hypothesis formulation and manuscript drafting, ensuring that analyses and writing are grounded in existing evidence. For research idea generation, another recent work~\cite{qi2023largelanguagemodelszero} constrains the process with curated background packets, using retrieval as an experimental control.


Building on these developments, retrieval-augmented generation (RAG) frameworks position external sources not merely as supporting references but as active components of the reasoning process. Specifically, RAG-enhanced scientific agents make external sources as primary inputs to LLM context and reasoning material, mostly with explicit planning, passage extraction, citation and contradiction checks. For example, PaperQA~\cite{lála2023paperqaretrievalaugmentedgenerativeagent} and PaperQA2~\cite{skarlinski2024languageagentsachievesuperhuman} treat retrieval as the main loop. By deciding which documents to read, attributing every claim, and detecting conflicts to steer synthesis, these works can yield expert-level literature reviews that are inherently verifiable. In material science, LLaMP~\cite{chiang2024llamp} extends RAG beyond text. Specifically, it utilizes  hierarchical ReAct~\cite{yao2023react} agents call material-specific APIs to fetch band gaps or elastic tensors, edit structures and then reason with computed properties.



\subsubsection{Self-evolving agentic reasoning}
\label{sec:app-science-evolving}

Scientific discovery agents can go beyond static reasoning and acquire the ability to self-evolve, which is to learn from experience, refine their internal representations and improve decision quality over successive interactions. This self-evolving layer equips agents with mechanisms to monitor and revise their own reasoning, retain and reuse intermediate hypotheses and adjust future plans based on external feedback or environmental signals. In the following paragraphs, we discuss how memory modules enable the accumulation  of scientific knowledge and how feedback and reflection mechanisms support continual adaptation and reasoning consistency throughout long-horizon scientific workflows.


\paragraph{Memory.}
\label{sec:app-science-memory}

ChemAgent~\cite{tang2025chemagentselfupdatinglibrarylarge} implements a self-updating library. It decomposes chemistry problems into sub-tasks and writes reusable \textit{skills} (ex: procedures, patterns, solutions) that later prompts can retrieve and adapt, stabilizing long multi-step reasoning without re-deriving everything from scratch. On the other hand, MatAgent~\cite{takahara2025acceleratedinorganicmaterialsdesign} emphasizes interpretable generation for inorganic materials, where \textit{short-term memory} recalls recent compositions and feedback, \textit{long-term memory} preserves successful designs together with their reasoning traces, and both are reused across iterations to guide proposal refinement and enable transparent audit.

\paragraph{Agentic Feedback and Reflection.}
\label{sec:app-science-feedback-reflect}


Firstly, Scientific Generative Agent~\cite{ma2024llmsimulationbileveloptimizers} ties discrete LLM proposals to inner-loop simulations that optimize continuous parameters, advancing only when evidence improves. The reflection ability is driven by measurable loss reductions. Next, ChemReasoner~\cite{sprueill2024chemreasoner} performs heuristic search over the LLM’s idea space but scores and steers candidates with quantum-chemical feedback, turning electronic-structure signals into a principled critique of linguistic hypotheses. Complementing these physics-based signals, Curie~\cite{kon2025curierigorousautomatedscientific} embeds rigor check directly into control flow via intra-agent checks, inter-agent gates and an experiment-knowledge module. In parallel, LLMatDesign~\cite{jia2024llmatdesignautonomousmaterialsdiscovery} builds explicit self-reflection into materials workflows, prompting the agent to surface and repair inconsistencies before they propagate to tool calls. Moreover, NovelSeek~\cite{novelseekteam2025novelseekagentscientist} utilizes reflection as a closed loop, updating code and plans with human-interactive feedback after each round. Finally, a recent study~\cite{kumbhar2025hypothesisgenerationmaterialsdiscovery} regularizes the process up front with explicit goals \& constraints and afterwards with standardized scoring to provides an objective standard that makes reflection repeatable.



\subsubsection{Collective multi-agent reasoning}
\label{sec:app-science-multi}



Multi-agent frameworks for scientific discovery distribute labor across specialized LLM-driven roles, where advanced LLM reasoning not only orchestrates coordination between scientific agents but also adjudicates conflicting evidence to maintain coherence in the process. 

To illustrate, we introduce some important multi-agent frameworks as follows. Firstly, ProtAgents~\cite{ghafarollahi2024protagents} exemplifies this pattern in protein design. The framework involve agents for literature retrieval, structure analysis, physics simulation, and results analysis. Specifically, the backbone LLM directs reasoning over multi-modal outputs, choosing when to iterate or convergence-check based on feedback signals. PiFlow~\cite{pu2025piflow}, on the other hand, instantiates reasoning as principle-aware uncertainty reduction with a multi-agent loop in which a Planner agent relays strategy to a Hypothesis agent and a validation loop, explicitly tying multi-agent communication to hypothesis–evidence alignment.  AtomAgents~\cite{ghafarollahi2024atomagentsalloydesigndiscovery} also brings similar role specialization to alloy discovery. In particular, the agent uses LLM-guided reasoning to control over when to trigger simulations and how to evaluate multi-modal results, letting reasoning allocate computational resources and prune alloy candidates.



With a similar \textit{planner, executor and evaluator} framework, CellAgent~\cite{xiao2024cellagentllmdrivenmultiagentframework} instantiates researches on single-cell analysis, where the planner LLM reasoning selects tools or hyper-parameters and the evaluator LLM triggers self-iterative re-runs when quality checks fail. Some other notable works include ARIA~\cite{ramirezmedina2025acceleratingscientificresearchmultillm} that introduces a four-agent framework (scout, filter, synthesizer and procedure-drafter), Curie~\cite{kon2025curierigorousautomatedscientific} that embeds rigor into multi-agent planning, Team of AI-made Scientists (TAIS)~\cite{liu2024toward} for gene-expression discovery and the Virtual Lab~\cite{swanson2024virtual} for nano-body design with role agents. 
